% Kapitel: Das hp-FEM-Plugin (Gitter mit adaptiven hp-Finiten-Elementen)
% Quellen: python/plugins/hpfem/ (hpfem_*.py, plugin.json, README.md), CHANGELOG.md,
% RTT-Fusion src/rtt_fusion/graph/hpfem_grating.py, adapters/hpfem.py (Stufe R35)
\chapter{Gitter mit hp-FEM}\label{kap:hpfem}
\index{hp-FEM}\index{Finite Elemente}\index{Beugungsgitter}

Das Plugin \plugin{hp-FEM} rechnet Beugungsgitter mit Finiten Elementen. Es verbindet \produkt{}
mit dem Rechenkern \emph{hp-FEM}: Nédélec- und H1-Elemente beliebiger Polynomordnung $p$ auf einem
Netz der Elementarzelle, periodische (Bloch-)Ränder links und rechts, absorbierende Schichten
(PML) oben und unten. Anders als die Fourier-Modal-Methode von FourierGrating
(\autoref{kap:fouriergrating}) zerlegt die Finite-Elemente-Methode die Zelle nicht in Schichten
mit konstantem Querschnitt: Schräge Flanken, Rundungen, Einschlüsse und Metalle mit scharfen
Kanten rechnet sie direkt auf dem Netz. Dafür kostet eine Rechnung mehr Speicher und Zeit, und
Netz und Ordnung $p$ bestimmen die Genauigkeit.

Das Plugin arbeitet mit den Aufträgen (Jobs) von hp-FEMs Job-Runner (\code{python -m hpfem.run},
Schema-Version 1, Problem \code{grating}). Alles, was einen Job ausmacht -- Zelle, Materialien,
Hintergrundstapel, Einfall, Netz und Löser --, steht im Panel \gui{Plugin}; ein gespeicherter Job
rechnet auch ohne \produkt{} mit dem Runner. Die Ergebnisse erscheinen als Tabellen, Balken,
Kurven und Bilder in den Ergebnisfenstern.

Das Kapitel führt durch die Zelle und ihre Koordinaten, Materialien und Stapel, Einfall, Netz
und Löser, dann durch die Rechnungen (Prüfen, Lösen, Feldkarten), die Zelle aus CAD-Körpern und
das Gitter als Block im Node-Workflow. Arbeitsabläufe, typische Meldungen und Grenzen schließen
es ab. Alle Parameter mit Standardwerten stehen im Anhang, \autoref{kap:plugin-hpfem}.

\section{Überblick}\label{sec:hp-ueberblick}

\subsection{Was das Plugin kann}

\begin{itemize}
  \item \textbf{Jobs laden und speichern} (\gui{Load job}, \gui{Save job}): das JSON-Format von
    hp-FEMs Job-Runner; Laden und erneutes Speichern ergibt dieselbe Datei.
  \item \textbf{Prüfen} (\gui{Check}): das Netz bauen und den Speicher einer Faktorisierung
    schätzen, ohne zu lösen.
  \item \textbf{Lösen} (\gui{Solve}): die Effizienz jeder reflektierten und transmittierten
    Ordnung, $\Sigma R$, $\Sigma T$, die Absorption $A$ und die Energiebilanz; für einen Sweep über
    die Wellenlänge oder den Einfallswinkel dieselben Größen als Kurven.
  \item \textbf{Feldkarten} (\gui{Field map}): $\vec E$, $\vec H$ oder den Poynting-Vektor
    $\vec S$ auf einem Raster über eine oder mehrere Perioden im Bildfenster.
  \item \textbf{Zelle zeigen} (\gui{Show cell}): die Tags der Zelle als Bild, ohne zu rechnen.
  \item \textbf{Zelle aus CAD-Körpern} (\gui{Shapes from CAD bodies}): im CAD gezeichnete Körper
    einer Elementarzelle als Polygone der Zelle.
  \item \textbf{Gitterblock im Node-Workflow}: \block{hpfem.grating} beugt die Strahlen eines
    optischen Aufbaus mit hp-FEM.
\end{itemize}

\subsection{Voraussetzungen}\label{sec:hp-voraussetzungen}
\index{hpfem (Python-Paket)}

Das Plugin wird mit der Anwendung installiert (Kennung \code{hpfem}, Name \plugin{hp-FEM}). Es
braucht in der Python-Umgebung der Plugins \code{numpy} und \code{hpfem} in Version 0.5 oder neuer
(\code{pip install hpfem}; Wheels gibt es für Linux x86\_64, Windows x86\_64 und macOS arm64 ab
macOS~15, sonst wird aus den Quellen gebaut). Fehlt \code{hpfem}, meldet jede rechnende Aktion:

\begin{lstlisting}
The hp-FEM plugin needs the Python package hpfem (>= 0.5) in the plugin environment
(pip install hpfem).
\end{lstlisting}

Laden und Speichern von Jobs, das Bearbeiten im Panel und die beiden Blockaktionen
funktionieren ohne das Paket. Drei Funktionen haben weitere Voraussetzungen:

\begin{itemize}
  \item Das Netz \gui{gmsh} braucht das Python-Paket \code{gmsh}.
  \item \gui{Shapes from CAD bodies} braucht Nurbforges Python-Paket \code{nurbforge} ab
    Version~0.7 (ebene Schnitte geschlossener Körper).
  \item Den Block \block{hpfem.grating} bietet das Plugin nur an, wenn \code{rtt-fusion} ab
    Version~0.5 installiert ist (Modul \code{rtt\_fusion.graph.hpfem\_grating}).
\end{itemize}

\subsection{Wo Sie das Plugin finden}

Wie jedes Plugin: im Menü \menu{Plugins,hp-FEM}, als eigene Gruppe in der Werkzeugspalte und im
Dock \gui{Plugin} (oben \gui{hp-FEM} wählen). Die zwei Aktionen für Gitterblöcke stehen nur im
Kontextmenü eines Blocks \block{hpfem.grating} im Node-Editor (\autoref{sec:hp-block}). Die
Parameter werden mit dem Dokument gespeichert.

\subsection{Aufbau des Panels}\label{sec:hp-panel}

Ganz oben steht ohne Rahmen die Liste \gui{Shapes (the last on top)}, darunter die Gruppen:

\begin{referenz}{>{\raggedright\arraybackslash}p{2.6cm}>{\raggedright\arraybackslash}X}
\toprule
Gruppe & Inhalt \\
\midrule
\endhead
\gui{File} & \gui{Job}, \gui{Save as}, \gui{Results in folder of} \\
\gui{Structure} & \gui{Name}, \gui{Unit}, \gui{Period}, \gui{Cell bottom y}, \gui{Cell top y},
  \gui{Background tag}, \gui{Slabs (full period)} \\
\gui{Materials} & \gui{Materials by tag} \\
\gui{Stack} & \gui{Incidence medium (above)}, \gui{Substrate (below)}, \gui{Background layers (top
  down)}, \gui{Top interface y} \\
\gui{Incidence} & \gui{Polarisation}, \gui{Wavelength}, \gui{$\theta$ / deg}, \gui{$\varphi$ / deg} \\
\gui{Mesh} & \gui{Mesh}, \gui{Columns (structured)}, \gui{Row bands (structured, bottom up)},
  \gui{Interface factor (gmsh)}, \gui{Mesh file (file)}, \gui{Scale to metres (file)} \\
\gui{Solver} & \gui{Polynomial order p}, \gui{PML top}, \gui{PML bottom}, \gui{Bottom},
  \gui{Orders up to |m|}, \gui{Diagnostics before solving} \\
\gui{Results} & \gui{Orders from efficiency}, \gui{List evanescent orders} \\
\gui{Field map} & \gui{Field}, \gui{Show}, \gui{Point of the sweep}, \gui{Scattered field only},
  \gui{Periods}, \gui{x from}/\gui{x to}, \gui{Points x}, \gui{y from}/\gui{y to}, \gui{Points y},
  \gui{Also write .npz} \\
\gui{CAD cell} & \gui{Bodies}, \gui{Bodies drawn in}, \gui{Section at CAD y}, \gui{Points per
  curved edge}, \gui{Shapes} \\
\gui{Advanced} & \gui{Further keys (JSON)} \\
\bottomrule
\end{referenz}

Der Standard ist das Beispiel des Runners: ein dielektrischer Steg ($\varepsilon = 2{,}25$,
200\,nm breit, 148\,nm hoch) auf SiO$_2$, Periode 400\,nm, p-polarisiert unter $\theta = 50°$,
$\varphi = 30°$ bei 400\,nm.

% ---------------------------------------------------------------------------------------------
\section{Die Zelle}\label{sec:hp-zelle}

\subsection{Koordinaten und Einheit}\label{sec:hp-koordinaten}
\index{hp-FEM!Koordinaten}

hp-FEM rechnet ein Liniengitter in einer zweidimensionalen Zelle:

\begin{itemize}
  \item $x$ läuft entlang der Periode, die Zelle ist \gui{Period} breit; ohne weitere Angabe reicht
    sie von $-\Lambda/2$ bis $\Lambda/2$ (ein anderer linker Rand: \code{model.x0} unter
    \gui{Advanced}).
  \item $y$ ist die Normale des Stapels; das Licht kommt von $+y$, also von oben. Die Zelle reicht
    von \gui{Cell bottom y} bis \gui{Cell top y}. Die PML-Schichten liegen \emph{innerhalb} dieser
    Grenzen.
  \item $z$ läuft entlang der Linien; in $z$ ist das Gitter unendlich und gleichförmig.
\end{itemize}

Alle Längen und die Wellenlänge stehen in der \gui{Unit} der Gruppe \gui{Structure} (nm, µm, mm
oder m).

\subsection{Tags, Schichten und Formen}
\index{Tag (hp-FEM)}

Jeder Punkt der Zelle hat einen \emph{Tag}, eine ganze Zahl, die über \gui{Materials by tag} sein
Material bestimmt. Welcher Tag gilt, entscheidet diese Reihenfolge:

\begin{enumerate}
  \item die letzte Form der Liste \gui{Shapes (the last on top)}, die den Punkt enthält -- Formen
    wiederholen sich periodisch;
  \item sonst die Schicht aus \gui{Slabs (full period)}, in deren $y$-Bereich (\gui{y0} bis
    \gui{y1}) er liegt -- Schichten reichen über die ganze Periode;
  \item sonst der \gui{Background tag}.
\end{enumerate}

Formen gibt es in fünf Typen: \code{rectangle} (linke untere Ecke \gui{x}, \gui{y}, \gui{width},
\gui{height}), \code{trapezoid} (Mitte \gui{x}, Basis \gui{y}, \gui{bottom}, \gui{top},
\gui{height}), \code{ellipse} (Mitte \gui{x}, \gui{y}, Halbachsen \gui{rx}, \gui{ry}),
\code{polygon} (\gui{Points (polygon)} gegen den Uhrzeigersinn) und \code{json} für jede andere
Form des Job-Formats. \gui{Show cell} zeigt das Ergebnis, bevor Sie rechnen
(\autoref{sec:hp-showcell}).

\subsection{Materialien}\label{sec:hp-materialien}
\index{hp-FEM!Materialien}

Ein Material schreiben Sie als eine Zeile Text:

\begin{referenz}{>{\raggedright\arraybackslash}p{3.6cm}>{\raggedright\arraybackslash}X}
\toprule
Eingabe & Bedeutung \\
\midrule
\endhead
\code{SiO2}, \code{Au}, \code{air} & ein Material aus hp-FEMs Bibliothek (dispersiv, für jede
  Wellenlänge ausgewertet; \gui{Library materials} listet sie im Log: Ag, Al, Al2O3, Au, Cr, GaAs,
  MAPbI3, Si, SiO2, Ti, TiO2, air, vacuum, water) \\
\code{1.5}, \code{3.88+0.02i} & Brechzahl $n$ bzw.\ $n + \mathrm i k$, unabhängig von der
  Wellenlänge \\
\code{eps=2.25}, \code{eps=-10+1.2i} & Permittivität $\varepsilon$ \\
\code{\{...\}} & jede andere Form des Job-Formats als JSON \\
\bottomrule
\end{referenz}

Bibliotheksmaterialien gelten nur in ihrem Wellenlängenbereich; außerhalb bricht die Rechnung mit
einer Meldung wie \enquote{TiO2: wavelength 0.4-0.4 um outside the range of validity 0.43-1.53 um}
ab.

\subsection{Der Hintergrundstapel}

Oberhalb und unterhalb der Formen sieht hp-FEM einen ebenen Schichtstapel: das
\gui{Incidence medium (above)}, die \gui{Background layers (top down)} mit ihren Dicken, das
\gui{Substrate (below)}; die oberste Grenzfläche liegt bei \gui{Top interface y}. Der Stapel
beschreibt die einfallende Welle und das Feld ohne Gitter; die Formen und Schichten der Zelle
sind die Störung darin. Er muss deshalb zu den Schichten der Zelle passen -- im Beispiel Luft über
SiO$_2$ mit der Grenzfläche bei $y = 0$, wo auch die SiO$_2$-Schicht der Zelle endet. Das
Einfallsmedium muss verlustfrei sein.

% ---------------------------------------------------------------------------------------------
\section{Einfall, Netz und Löser}\label{sec:hp-einfall}

\subsection{Einfall und Sweep}\label{sec:hp-sweep}

\gui{Polarisation} \code{s} ist bei $\varphi = 0$ das Feld $E_z$ entlang der Linien (TE),
\code{p} das Feld $H_z$ (TM). \gui{$\theta$ / deg} ist der Polarwinkel gegen die Normale,
\gui{$\varphi$ / deg} der Azimut; $\varphi \neq 0$ ist konischer Einfall.

\gui{Wavelength} und \gui{$\theta$ / deg} nehmen eine Zahl (\code{400}), eine Liste
(\code{400, 410, 425}) oder einen Bereich \code{Start:Stopp:Anzahl} (\code{400:700:31}). Eine
Liste oder ein Bereich ist ein \emph{Sweep}; Sie können die Wellenlänge oder $\theta$
durchfahren, nicht beide zugleich.

\subsection{Das Netz}\label{sec:hp-netz}
\index{hp-FEM!Netz}

\begin{description}
  \item[\code{structured}] Spalten mal Zeilenbänder, ohne Gmsh: \gui{Columns (structured)} gleich
    breite Spalten über die Periode und je Band in \gui{Row bands (structured, bottom up)} von
    \gui{y0} bis \gui{y1} \gui{rows} Zeilen. Jedes Element bekommt den Tag an seiner Mitte; legen
    Sie Grenzflächen und Kanten der Formen auf Gitterlinien. Für Rechtecke ist das exakt und
    schnell.
  \item[\code{gmsh}] ein Dreiecksnetz, dessen Elementgrößen hp-FEM aus der Wellenlänge im Material
    und der Ordnung $p$ wählt; an Grenzflächen verfeinert \gui{Interface factor (gmsh)}. Richtig
    für schräge und runde Formen. Braucht das Paket \code{gmsh}.
  \item[\code{file}] ein periodisches Gmsh-Netz (\datei{.msh}) aus einer Datei; \gui{Scale to
    metres (file)} rechnet seine Koordinaten in Meter um. Ein relativer Pfad gilt relativ zur
    Job-Datei.
\end{description}

\subsection{Der Löser}\label{sec:hp-loeser}

\begin{description}
  \item[\gui{Polynomial order p}] Ordnung der Elemente; höher ist genauer und teurer. Sind die
    Elemente für die Wellenlänge im Material zu grob, sagt es die Diagnose im Log.
  \item[\gui{PML top}, \gui{PML bottom}] Dicke der absorbierenden Schichten innerhalb der Zelle;
    beide 0: hp-FEMs Empfehlung.
  \item[\gui{Bottom}] \code{pml} (offenes Substrat), \code{pec} (Metallwand, keine transmittierten
    Ordnungen) oder \code{modal} (modaler Abschluss, für Substrate, in denen eine Ordnung rückwärts
    läuft).
  \item[\gui{Orders up to |m|}] die Ordnungen $-m \ldots m$, die ausgewertet werden.
  \item[\gui{Diagnostics before solving}] hp-FEMs Prüfungen (Elemente je Wellenlänge, PML, Lage der
    Auswertelinien) mit Hinweisen im Log.
\end{description}

\begin{hinweis}
Die Genauigkeit hängt an Netz und Ordnung. Rechnen Sie einen Punkt mit $p$ und $p+1$ (oder einem
feineren Netz): Ändern sich die Effizienzen kaum, ist die Lösung konvergiert. Die Energiebilanz
(\code{balance}) eines verlustfreien Gitters ist ein zweiter Anhalt.
\end{hinweis}

% ---------------------------------------------------------------------------------------------
\section{Rechnen}\label{sec:hp-rechnen}

\subsection{Prüfen: Check}\label{sec:hp-check}

\gui{Check} baut das Netz und schätzt den Speicher einer Faktorisierung, ohne zu lösen. Im Log
stehen das Netz (Zellen, Knoten, Tags ohne Material, Materialien ohne Zellen, ungültige Zellen)
und die Schätzung; braucht eine Faktorisierung mehr als 70\,\% des Arbeitsspeichers, warnt eine
Zeile \enquote{warning: one factorisation needs about … GiB}. Die Statusleiste meldet die Zahl der
Zellen und Punkte und \enquote{ready to solve} -- oder, was im Log zu prüfen ist.

\subsection{Lösen: Solve}\label{sec:hp-solve}

\gui{Solve} rechnet den Job. Der Fortschrittsbalken zeigt Punkt und Phase (Assemblieren,
Faktorisieren, Lösen); \gui{Cancel} hält zwischen den Phasen an. Je Punkt steht eine Zeile im Log.

\begin{itemize}
  \item \textbf{Ein Punkt}: Ergebnissatz \enquote{\textit{Name}: $\lambda$ = …, $\theta$ = …°} mit der Tabelle
    \gui{Orders} (Spalten \code{side}, \code{m}, \code{efficiency}, \code{propagating}) und den
    Balken \gui{Efficiency per order}; die Statusleiste zeigt $\Sigma R$, $\Sigma T$, $A$ und die
    Bilanz. \gui{List evanescent orders} nimmt die abklingenden Ordnungen in die Tabelle auf.
  \item \textbf{Ein Sweep}: die Kurven $\Sigma R$, $\Sigma T$, $A$ und jede Ordnung
    (\enquote{R~1}, \enquote{T~0}, …), die irgendwo \gui{Orders from efficiency} erreicht, über
    $\lambda$ oder $\theta$, dazu die Tabelle \gui{Points} (Punkt, $\lambda$, $\theta$, R, T, A,
    Bilanz, Freiheitsgrade, Sekunden).
\end{itemize}

Die Datei \datei{results.json} des Runners landet im Ordner von \gui{Results in folder of} oder
in einem temporären Ordner; der Pfad steht im Log.

\subsection{Feldkarten: Field map}\label{sec:hp-feldkarte}
\index{hp-FEM!Feldkarte}

\gui{Field map} rechnet einen Punkt des Sweeps (\gui{Point of the sweep}, von 1 an gezählt) und
zeigt ein Feld auf einem Raster im Fenster \gui{Image}:

\begin{itemize}
  \item \gui{Field}: \code{E} (einfallende Welle mit 1\,V/m), \code{H} in A/m oder der
    Poynting-Vektor \code{S} in W/m²; mit \gui{Scattered field only} nur das Streufeld.
  \item \gui{Show}: $|F|^2$, die Beträge $|F_x|$, $|F_y|$, $|F_z|$ oder Real- und Imaginärteil einer
    Komponente (Farbskala symmetrisch um 0). $\mathrm{Re}\,S_y < 0$ ist Leistung, die nach unten
    fließt.
  \item Das Raster: in $x$ über \gui{Periods} Perioden vom linken Zellrand (oder von \gui{x from}
    bis \gui{x to}), in $y$ über die Zelle ohne die PML-Schichten des Lösers (oder von \gui{y from}
    bis \gui{y to}); \gui{Points x}, \gui{Points y} Punkte. Punkte außerhalb der Zelle bleiben leer.
  \item \gui{Also write .npz} speichert die Karte des Runners (\code{x}, \code{y} in Metern,
    \code{values} der Form $(n_y, n_x, 3)$).
\end{itemize}

\subsection{Die Zelle ansehen: Show cell}\label{sec:hp-showcell}

\gui{Show cell} zeichnet die Tags der Zelle auf demselben Raster wie die Feldkarte, ohne zu
rechnen, und schreibt das Material jedes Tags ins Log. Formen, die es nicht zeichnen kann
(JSON-Formen eines unbekannten Typs), nennt eine Warnung.

% ---------------------------------------------------------------------------------------------
\section{Jobs als Dateien}\label{sec:hp-dateien}

\gui{Load job} liest die Datei aus \gui{Job} ins Panel, \gui{Save job} schreibt das Panel nach
\gui{Save as} (oder, wenn leer, nach \gui{Job}). Was das Panel nicht zeigt -- zum Beispiel
Feldkarten des Runners, \code{model.x0} oder Toleranzen des Lösers --, steht unter \gui{Advanced}
als JSON und wird beim Speichern wieder eingefügt. Ein gespeicherter Job rechnet mit

\begin{lstlisting}
python -m hpfem.run job.json --out ergebnisse
\end{lstlisting}

% ---------------------------------------------------------------------------------------------
\section{Die Zelle aus CAD-Körpern}\label{sec:hp-cad}
\index{hp-FEM!CAD-Körper}

Eine Zelle mit schrägen Flanken, Rundungen oder Löchern zeichnen Sie bequemer im CAD als in der
Liste der Formen. \gui{Shapes from CAD bodies} schneidet die gewählten Körper quer zu den Linien
und macht aus jedem Schnitt Polygone der Zelle.

\subsection{Die Körper zeichnen}

\begin{itemize}
  \item Die Linien laufen entlang CAD~$+Y$, die Normale des Stapels entlang CAD~$+Z$ (Licht von
    oben). Das $x$ des Jobs ist CAD~$x$, sein $y$ ist CAD~$z$.
  \item Zeichnen Sie geschlossene Körper (Schalen) einer Periode, in $y$ lang genug, dass die
    Schnittebene sie trifft.
  \item \gui{Bodies drawn in} sagt, was eine CAD-Längeneinheit ist (nm, µm oder mm); sie wird in
    die \gui{Unit} des Jobs umgerechnet.
  \item Das Material eines Körpers ist sein Attribut \code{material} (jede Form aus
    \autoref{sec:hp-materialien}) oder, ohne Attribut, sein Name, wenn er ein Material des Panels
    oder der Bibliothek ist.
\end{itemize}

\subsection{Die Formen anlegen}

Wählen Sie die Körper im Modell, übernehmen Sie sie in \gui{Bodies} und starten Sie
\gui{Shapes from CAD bodies}:

\begin{itemize}
  \item Nurbforge schneidet jeden Körper exakt mit der Ebene CAD $y$ = \gui{Section at CAD y}.
    Gerade Kanten des Schnitts behalten ihre Endpunkte, gekrümmte bekommen \gui{Points per curved
    edge} Punkte.
  \item Jede Schleife wird ein Polygon gegen den Uhrzeigersinn. Eine Schleife innerhalb einer
    anderen desselben Körpers ist ein Loch: Es wird ein Polygon obendrauf mit dem Tag, der dort
    ohne den Körper gilt (Schicht oder Hintergrund).
  \item Körper desselben Materials teilen sich einen Tag (verglichen wird der Wert:
    \code{eps=2.25} ist dasselbe wie \code{eps=2.25+0i}); ein neues Material bekommt den nächsten
    freien Tag in \gui{Materials by tag}.
  \item \gui{Shapes} ersetzt die Formen des Panels (\code{replace shapes}) oder legt die Körper
    obendrauf (\code{add shapes}).
\end{itemize}

Das Log nennt den Bereich der Körper, Körper, die die Ebene verfehlen, und Körper, die die Zelle
verlassen -- Periode, Zellgrenzen und Stapel bleiben die des Panels. Mit dem strukturierten Netz
bekommt jedes Element den Tag an seiner Mitte; für schräge oder runde Körper wählen Sie
\gui{Mesh} \code{gmsh}.

% ---------------------------------------------------------------------------------------------
\section{Das Gitter als Block im Workflow}\label{sec:hp-block}
\index{hpfem.grating@\texttt{hpfem.grating} (Block)}

Im Node-Workflow (\autoref{kap:workflow}) liefert das Plugin den Block \block{hpfem.grating}: eine
periodische Probe, die jede gewünschte Beugungsordnung als eigenen Kanal weitergibt -- wie
\block{fg.grating} von FourierGrating (\autoref{sec:fg-block}), aber mit hp-FEM gerechnet. Er hat
die Anschlüsse eines Spiegels (\code{in}, \code{out} \enquote{orders}); \gui{Replace with} (\autoref{sec:workflow-ersetzen})
behält deshalb die Verbindungen. Die Workflow-Engine schickt seine Auswertung an den Prozess des
hp-FEM-Plugins.

\subsection{Der Job des Blocks}

Das Gitter des Blocks ist ein hp-FEM-Job ohne Einfall, Sweep und Feldkarten; den Einfall setzen
die Strahlen, die am Block ankommen. Der Block hält den Job selbst (Parameter \param{job\_data}),
ein \datei{.rttf}-Dokument rechnet damit auf jedem Rechner; eine Netzdatei (\gui{Mesh}
\code{file}) bleibt allerdings außerhalb des Dokuments. Zwischen Panel und Block wechseln Sie im
\textbf{Kontextmenü des Blocks} im Node-Editor:

\begin{description}
  \item[\gui{Store job in block}] Der Job des Panels wird der des Blocks. Einfall, Sweep und Karten
    fallen weg (das Log sagt es), ein relativer Netzpfad wird absolut. Die Statusleiste meldet die
    Zahl der Materialien und Formen.
  \item[\gui{Load block job into panel}] zeigt den Job des Blocks im Panel zum Bearbeiten, Prüfen
    oder Lösen; Einfall und Sweep bleiben die des Panels. \gui{Job} wird geleert, damit
    \gui{Save job} keine fremde Datei überschreibt.
\end{description}

Im Dialog des Blocks stehen danach die wichtigsten Zahlen des Jobs als Felder: \gui{Period}
(Gruppe \gui{Grating}), \gui{Polynomial order} (\gui{Solver}) und je Material, das als $n$ oder
$n, k$ gegeben ist, \gui{Tag $t$: n} bzw.\ \gui{Tag $t$: n, k} (\gui{Materials}).

\subsection{Lage und Ordnungen}

Das $x$ des Jobs ist der \param{grating\_vector} des Blocks, sein $z$ die Linien, sein $y$ die
\param{normal}. Das $x = 0$ des Jobs liegt im \param{point} des Blocks, und dieser Punkt ist die
Oberkante der Struktur: der höchste Punkt der Formen und des Stapels. \param{side} wählt die
reflektierten oder die transmittierten Ordnungen, \param{orders} welche (leer: $-1$, 0, 1).

\subsection{Wie der Block rechnet}

Um die Strahlen einer Wellenlänge baut der Block eine Tabelle der Jones-Matrizen jeder Ordnung
(wie \block{fg.grating}, \autoref{sec:fg-block}):

\begin{itemize}
  \item Kollimiertes Licht (eine Richtung) rechnet er genau dort: eine hp-FEM-Rechnung mit zwei
    Lösungen je Richtung.
  \item Für einen Kegel von Richtungen legt er Stützstellen im Abstand \gui{$\sigma$ step of the
    table} über den Bereich der Richtungskosinus $\sigma = n \sin\theta$ plus \gui{$\sigma$ margin
    of the table}; je Stützstelle zwei Lösungen. Würden es mehr als \gui{Most table nodes} (400),
    bricht der Block mit einer Meldung ab, bevor hp-FEM rechnet.
  \item Tabellen und vorbereitete Zellen bleiben erhalten, solange das Plugin läuft; ein zweiter
    Lauf mit demselben Licht rechnet nicht neu.
\end{itemize}

Der Block rechnet in Strahlen; Felder und Spektren, die ihn erreichen, wandelt die Workflow-Engine
vorher in Strahlen um.

% ---------------------------------------------------------------------------------------------
\section{Arbeitsabläufe Schritt für Schritt}\label{sec:hp-ablaeufe}

\subsection{Die erste Rechnung}

\begin{enumerate}
  \item Wählen Sie im Dock \gui{Plugin} \gui{hp-FEM}. Das Panel zeigt das Beispiel des Runners.
  \item \gui{Show cell}: Das Bild zeigt den Steg (Tag 3) auf SiO$_2$ (Tag 2) unter Luft (Tag 1).
  \item \gui{Check}: Im Log stehen Netz und Speicherschätzung; die Statusleiste meldet
    \enquote{ready to solve}.
  \item \gui{Solve}: Die Tabelle \gui{Orders} zeigt die Effizienzen; $\Sigma R + \Sigma T$ ist
    bis auf die Bilanz 1.
\end{enumerate}

\subsection{Ein Spektrum}

Setzen Sie \gui{Wavelength} auf \code{400:700:31} und starten Sie \gui{Solve}. Die Kurven zeigen
$\Sigma R$, $\Sigma T$ und jede Ordnung, die 1\,\% erreicht, über $\lambda$; die Tabelle
\gui{Points} hat eine Zeile je Wellenlänge. Mit \gui{Cancel} behalten Sie die bis dahin
gerechneten Punkte.

\subsection{Das Feld am Steg}

Setzen Sie \gui{Point of the sweep} auf einen Punkt des Spektrums, \gui{Field} \code{E},
\gui{Show} $|F|^2$ und starten Sie \gui{Field map}. Mit \gui{Field} \code{S} und \gui{Show}
\code{Re Fy} sehen Sie, wohin die Leistung fließt.

\subsection{Eine Zelle aus dem CAD}

\begin{enumerate}
  \item Zeichnen Sie einen Körper einer Periode mit den Linien entlang CAD~$+Y$ -- etwa einen
    Zylinder mit der Achse in $Y$ als runden Draht -- und geben Sie ihm das Attribut
    \code{material} (z.\,B.\ \code{Au}).
  \item Wählen Sie ihn, übernehmen Sie ihn in \gui{CAD cell} $\to$ \gui{Bodies}, setzen Sie
    \gui{Bodies drawn in} und \gui{Section at CAD y}.
  \item \gui{Shapes from CAD bodies}, dann \gui{Show cell}. Stellen Sie \gui{Mesh} auf \code{gmsh}
    und rechnen Sie.
\end{enumerate}

\subsection{Ein hp-FEM-Gitter im Linnik}

\begin{enumerate}
  \item Ersetzen Sie im Node-Workflow den Probenspiegel durch einen Block \block{hpfem.grating}
    (\gui{Replace with}).
  \item Stellen Sie das Gitter im Panel ein und prüfen Sie es mit \gui{Solve}.
  \item \maus{Rechtsklick} auf den Block $\to$ \gui{Store job in block}.
  \item Setzen Sie im Dialog des Blocks \param{side} \code{reflected} und die gewünschten
    \param{orders}; \gui{Run workflow}.
\end{enumerate}

% ---------------------------------------------------------------------------------------------
\section{Typische Meldungen}\label{sec:hp-fehler}

\begin{description}
  \item[\enquote{The panel does not give a valid job: …}] Eine Eingabe passt nicht, z.\,B.\ ein
    Polygon mit weniger als drei Punkten, Wellenlänge und $\theta$ zugleich als Sweep oder ein
    Material wie \code{eps=x}.
  \item[\enquote{hp-FEM: …}] eine Meldung des Rechenkerns, etwa ein unbekanntes Material, eine
    Wellenlänge außerhalb des Gültigkeitsbereichs eines Bibliotheksmaterials oder PML-Dicken, die
    nicht in die Zelle passen.
  \item[\enquote{tags without material} im Log] Ein Tag der Zelle hat keine Zeile in
    \gui{Materials by tag}.
  \item[\enquote{warning: one factorisation needs about … GiB}] Die Rechnung passt kaum in den
    Speicher: Ordnung oder Netz verkleinern.
  \item[\enquote{body … has no material}] Ein CAD-Körper hat weder das Attribut \code{material}
    noch einen Namen, der ein Material ist.
  \item[\enquote{… table nodes (2 hp-FEM solves each) are more than …}] Das Strahlenbündel am
    Gitterblock ist zu breit für die Tabelle: \gui{$\sigma$ step of the table} oder \gui{Most table
    nodes} erhöhen.
\end{description}

\section{Grenzen}\label{sec:hp-grenzen}

\begin{itemize}
  \item Das Plugin rechnet Liniengitter (2D-Zelle, auch konischer Einfall), keine Kreuzgitter.
  \item Das Panel zeigt Jobs der Aufgabe \code{scattering}; andere Aufgaben des Runners (Resonanzen,
    Bänder, Emitter, Studien) laufen nur mit dem Runner selbst.
  \item Konvergenz in Netz und Ordnung prüft das Plugin nicht; das ist Sache der Rechnung
    (\autoref{sec:hp-loeser}).
  \item Der Gitterblock rechnet Strahlen einer Wellenlänge je Tabelle; ein Feld oder Spektrum wird
    vorher zu Strahlen.
\end{itemize}
